\chapter{Metodología}
\label{chap:metodologia}

Este capítulo sostiene que una metodología de gestión de riesgos sobre datos personales solo es válida si parte de la práctica real de la institución, se construye con un método de investigación verificable y traduce cada exigencia normativa en una actividad concreta. Para demostrarlo, se caracteriza primero la situación de partida de la Universidad Tecnológica Indoamérica (UTI) en el proceso de admisión y matrícula de la sede Quito; se justifica después el enfoque de ciencia del diseño y las técnicas de recolección; y se presentan, por último, los requisitos y el diseño de la metodología propuesta. La información institucional se obtiene con autorización formal de la universidad y se trata de forma anonimizada, de acuerdo con las limitaciones declaradas en el Capítulo~\ref{chap:introduccion}.

\section{Diagnóstico de la Situación Actual}
\label{sec:diagnostico_actual}

Un diagnóstico honesto no puede presuponer las deficiencias de la institución que analiza. Por ello, esta sección distingue con claridad entre lo que se conoce de la UTI por fuentes públicas verificables, lo que se sabe del proceso de admisión y matrícula por la literatura y la práctica del sector, y lo que solo podrá establecerse con la aplicación de los instrumentos de la fase de diagnóstico (OE1), cuyos resultados se presentan en el Capítulo~\ref{chap:implementacion}.

\subsection{Caracterización de la institución}

La Tabla~\ref{tab:parametros_ies} reúne la información pública verificada de la UTI y señala los datos que se levantarán con autorización institucional. Los valores marcados como pendientes no se estiman ni se suponen, porque de ellos dependen criterios de impacto exigidos por la guía de la SPDP, como el número de titulares afectados y la presencia de categorías especiales \cite{spdp2026guiav2}.

\begin{table}[htbp]
  \centering
  \caption{Caracterización de la Universidad Tecnológica Indoamérica para el caso de estudio.}
  \label{tab:parametros_ies}
  \footnotesize
  \renewcommand{\arraystretch}{1.25}
  \begin{tabularx}{\textwidth}{|p{4.4cm}|L|p{3.0cm}|}
    \hline
    \textbf{Parámetro} & \textbf{Valor} & \textbf{Fuente} \\
    \hline
    Creación & Ley N.º 112, Registro Oficial N.º 373, 31 de julio de 1998 & \cite{congreso1998uti} \\
    \hline
    Naturaleza jurídica & Entidad de derecho privado, con personería jurídica y autonomía administrativa y financiera & \cite{congreso1998uti} \\
    \hline
    Domicilio principal & Ambato, provincia de Tungurahua & \cite{congreso1998uti} \\
    \hline
    Financiamiento & Recursos propios, matrículas, pensiones y aranceles, legados y donaciones & \cite{congreso1998uti} \\
    \hline
    Sede del caso de estudio & Quito & Información institucional \\
    \hline
    Estudiantes matriculados en la sede Quito & \textcolor{red}{[PENDIENTE]} & Secretaría Académica \\
    \hline
    Aspirantes por periodo y proporción de menores de edad & \textcolor{red}{[PENDIENTE]} & Admisiones \\
    \hline
    Estudiantes con beca, arancel diferenciado o discapacidad declarada & \textcolor{red}{[PENDIENTE]} & Bienestar Estudiantil \\
    \hline
    Unidades que intervienen en el proceso & \pend{validar el modelo de la Tabla~\ref{tab:subprocesos}} & Entrevistas \\
    \hline
    Sistemas de información y encargados del tratamiento & \textcolor{red}{[PENDIENTE]} & Dirección de TI \\
    \hline
    Política de protección de datos y delegado de protección de datos & \pend{no se localizaron en la búsqueda pública realizada en octubre de 2026} & Área jurídica \\
    \hline
    Autorización formal de la universidad para el estudio & \pend{oficio de autorización, instancia que lo emite y fecha (Anexo~\ref{anx:autorizacion})} & Rectorado o instancia que la UTI designe \\
    \hline
  \end{tabularx}
  \raggedright
  \elaboradopor{Ramón Alejandro Manzaba Castro}
  \fuente{Elaboración propia con base en \cite{congreso1998uti}}
\end{table}

\subsection{Descripción del proceso de admisión y matrícula}

El proceso de admisión y matrícula se eligió porque concentra, en un periodo breve, la recolección masiva de datos de personas que aún no tienen vínculo con la institución, incluye previsiblemente categorías especiales de datos ---discapacidad, condición socioeconómica y datos de adolescentes--- e involucra a varias unidades y proveedores externos. Esta combinación lo vuelve representativo de los riesgos que la LOPDP pretende controlar \cite[arts.~25 y 40]{lopdp2021} y lo sitúa potencialmente dentro de los supuestos de tratamiento a gran escala definidos por la SPDP, que incluyen el tratamiento sistemático de datos de niñas, niños y adolescentes en entornos educativos \cite{spdp2026granescala}.

La Tabla~\ref{tab:subprocesos} propone un modelo de referencia del proceso en seis subprocesos, construido a partir de la estructura habitual de los procesos de admisión en las IES. El modelo no describe todavía la práctica de la UTI: es el punto de partida de las entrevistas, en las que cada responsable confirmará, corregirá o completará los datos, los sistemas y los terceros que intervienen en su subproceso.

\begin{table}[htbp]
  \centering
  \caption{Modelo de referencia del proceso de admisión y matrícula (a validar en la UTI).}
  \label{tab:subprocesos}
  \footnotesize
  \renewcommand{\arraystretch}{1.2}
  \begin{tabularx}{\textwidth}{|p{2.6cm}|L|p{2.4cm}|L|}
    \hline
    \textbf{Subproceso} & \textbf{Datos personales previsibles} & \textbf{Titulares} & \textbf{Sistemas y terceros previsibles} \\
    \hline
    SP1. Registro y postulación & Identificación, contacto, fecha de nacimiento, colegio de procedencia, fotografía & Aspirantes; representantes de menores & Portal web; herramienta de gestión de postulantes \\
    \hline
    SP2. Admisión y nivelación & Calificaciones de bachillerato, resultados de evaluaciones de admisión & Aspirantes & Sistema de gestión académica; correo institucional \\
    \hline
    SP3. Becas y aranceles diferenciados & Información socioeconómica del hogar, certificados de discapacidad, datos de salud en casos de excepción & Aspirantes y su núcleo familiar & Formularios y almacenamiento institucional \\
    \hline
    SP4. Pago de matrícula & Datos de facturación y de medios de pago del pagador & Estudiantes; pagadores & Sistema financiero; proveedor de pagos \\
    \hline
    SP5. Legalización de la matrícula & Expediente completo, creación de cuentas institucionales & Estudiantes & Sistema de gestión académica; servicios en la nube; aula virtual \\
    \hline
    SP6. Reportes y archivo & Datos para reportes a organismos de control; expedientes de aspirantes no admitidos & Estudiantes y aspirantes & Sistema de gestión académica; archivo \\
    \hline
  \end{tabularx}
  \raggedright
  \elaboradopor{Ramón Alejandro Manzaba Castro}
  \fuente{Elaboración propia}
\end{table}

\subsection{Aspectos críticos a verificar}

La literatura permite anticipar dónde suelen concentrarse los riesgos en procesos de este tipo, aunque no permite afirmar que existan en una institución concreta. Las revisiones de Ulven y Wangen \cite{ulven2021systematic} y de Álvarez-Carrión y Hernández-Sotomayor \cite{alvarez2024proteccion} señalan como causas recurrentes de incidentes en universidades el error humano, los fallos técnicos y la gestión deficiente de accesos, mientras que Li et al. \cite{li2023breach} vinculan la ocurrencia de brechas con decisiones organizacionales sobre vulnerabilidades y plataformas, y Piray Rodríguez et al. \cite{piray2026garantias} documentan una adopción desigual de políticas de protección de datos en las universidades del país. Sobre esa base, el diagnóstico de la UTI verificará, con evidencia de la práctica efectiva y no solo de la política escrita, los siguientes aspectos:

\begin{enumerate}
  \item La existencia y actualización de un registro de actividades de tratamiento (RAT) del proceso.
  \item La base de legitimación de cada actividad y la separación, en los formularios, de las finalidades necesarias para la admisión y las finalidades de difusión.
  \item El modo de recepción, almacenamiento y acceso de los certificados de discapacidad y la información socioeconómica.
  \item La asignación y revisión de los perfiles de acceso a los sistemas, en especial del personal temporal.
  \item Los plazos de conservación y eliminación de los expedientes, incluidos los de aspirantes no admitidos.
  \item Las cláusulas de protección de datos en los contratos con los encargados del tratamiento.
  \item La existencia de un procedimiento de gestión y notificación de vulneraciones que permita cumplir el término de cinco días del artículo~43 de la LOPDP \cite[art.~43]{lopdp2021}.
  \item El grado en que la gestión de la seguridad considera el impacto de los incidentes sobre los titulares.
\end{enumerate}

En consecuencia, el diagnóstico no parte de suponer deficiencias en la UTI, sino de una lista explícita de aspectos que la evidencia confirmará o descartará. Esta decisión protege la validez del trabajo y evita atribuir a la institución hallazgos que no se han comprobado.

\subsection{Requisitos para ejecutar el diagnóstico}
\label{sec:requisitos_diagnostico}

El primer objetivo específico exige saber con precisión qué se necesita para diagnosticar, antes de diagnosticar. La Tabla~\ref{tab:requisitos_diagnostico} enumera los requisitos de información, su fuente o informante, el instrumento con que se obtienen, su fundamento y su estado a la fecha de corte del documento. El detalle operativo ---qué documento solicitar, a quién y con qué formato--- se mantiene en el archivo de control \texttt{REQUISITOS\_DIAGNOSTICO.md} del repositorio del proyecto.

\begin{table}[htbp]
  \centering
  \caption{Requisitos para ejecutar el diagnóstico del proceso de admisión y matrícula (OE1).}
  \label{tab:requisitos_diagnostico}
  \scriptsize
  \renewcommand{\arraystretch}{1.2}
  \begin{tabularx}{\textwidth}{|c|L|>{\raggedright\arraybackslash}p{2.4cm}|>{\raggedright\arraybackslash}p{2.2cm}|>{\raggedright\arraybackslash}p{2.2cm}|>{\raggedright\arraybackslash}p{1.5cm}|}
    \hline
    \textbf{N.º} & \textbf{Requisito} & \textbf{Fuente o informante} & \textbf{Instrumento} & \textbf{Fundamento} & \textbf{Estado} \\
    \hline
    D01 & Autorización formal de la universidad y designación de un enlace institucional & Rectorado o instancia designada & Oficio (Anexo~\ref{anx:autorizacion}) & Ética del estudio; LOPDP, art.~7 & \pend{oficio} \\
    \hline
    D02 & Requisitos normativos aplicables al proceso & LOPDP, RLOPDP, resoluciones de la SPDP & Lista de verificación (Anexo~\ref{anx:lista_verificacion}) & LOPDP, arts.~37 y 40 & Disponible \\
    \hline
    D03 & Inventario de actividades de tratamiento y flujos de datos de SP1 a SP6: finalidades, categorías, titulares, destinatarios, encargados, transferencias y conservación & Responsables de Admisiones, Secretaría, Bienestar, Financiero y TI & Guía de entrevista (Anexo~\ref{anx:entrevista}); formato RAT (Tabla~\ref{tab:rat}) & RLOPDP, arts.~38 y 39 & \pend{entrevistas} \\
    \hline
    D04 & Volumen de titulares por periodo y proporción de aspirantes menores de edad & Admisiones; Secretaría Académica & Solicitud de datos agregados & Res.~SPDP-SPD-2026-0005-R (gran escala) & \pend{datos agregados} \\
    \hline
    D05 & Documentos vigentes: política y avisos de privacidad, formularios, contratos con encargados, tabla de conservación, matriz de accesos, procedimiento de incidentes & Área jurídica; Secretaría; TI & Lista de verificación (Anexo~\ref{anx:lista_verificacion}) & LOPDP, arts.~34, 37, 43 y 47 & \pend{documentos} \\
    \hline
    D06 & Controles técnicos y organizacionales existentes y evidencia de su práctica efectiva & Dirección de TI & Auditoría de procesos (Anexo~\ref{anx:lista_verificacion}) & LOPDP, art.~41; guía SPDP & \pend{evidencia} \\
    \hline
    D07 & Nivel de madurez por dominio & Responsables de las unidades & Cuestionario de madurez (Anexo~\ref{anx:madurez}) & Fase 0 de la metodología & \pend{respuestas} \\
    \hline
    D08 & Historial agregado de incidentes de los últimos tres años & Dirección de TI & Entrevista; registro de incidentes & LOPDP, art.~41, num.~4 & \pend{historial} \\
    \hline
    D09 & Existencia de delegado de protección de datos y su registro ante la SPDP & Área jurídica & Entrevista; consulta documental & LOPDP, art.~48 & \pend{confirmación} \\
    \hline
    D10 & Panel de al menos cinco expertos con coeficiente de competencia determinado & Profesionales externos al caso & Autoevaluación de competencia (Anexo~\ref{anx:competencia}) & Guía SPDP (Delphi) & \pend{panel} \\
    \hline
  \end{tabularx}
  \raggedright
  \elaboradopor{Ramón Alejandro Manzaba Castro}
  \fuente{Elaboración propia con base en \cite{lopdp2021,rlopdp2023,spdp2026guiav2,spdp2026granescala}}
\end{table}

Los requisitos D01 a D10 no son equivalentes: D01 condiciona a todos los demás, porque sin autorización no puede accederse a la información institucional, y D03 y D04 condicionan la determinación de gran escala y, con ella, la obligatoriedad de la evaluación de impacto. Por ese motivo, el protocolo de aplicación del Capítulo~\ref{chap:implementacion} los ordena en ese mismo sentido.

\section{Área de Estudio}

La Tabla~\ref{tab:area_estudio} sintetiza el dominio, el área y el objeto de estudio del proyecto.

\begin{table}[htbp]
  \centering
  \caption{Área de estudio.}
  \label{tab:area_estudio}
  \renewcommand{\arraystretch}{1.35}
  \begin{tabular}{|p{0.39\textwidth}|p{0.53\textwidth}|}
    \hline
    \textbf{Dominio:} & Gobernanza de la información y protección de datos personales \\
    \hline
    \textbf{Línea de investigación:} & \UTILineaInvestigacion \\
    \hline
    \textbf{Sublínea de investigación:} & \UTISublineaInvestigacion \\
    \hline
    \textbf{Campo:} & Tecnologías de la información y la comunicación (TIC) \\
    \hline
    \textbf{Área:} & Gestión de riesgos sobre el tratamiento de datos personales y evaluación de impacto \\
    \hline
    \textbf{Aspectos:} & Inventario de tratamientos, identificación de amenazas y vulnerabilidades, estimación de probabilidad e impacto sobre los titulares, tratamiento del riesgo, monitoreo y cumplimiento de la LOPDP, su reglamento y la normativa técnica de la SPDP \\
    \hline
    \textbf{Objeto de estudio:} & Proceso de admisión y matrícula de la Universidad Tecnológica Indoamérica, sede Quito \\
    \hline
    \textbf{Periodo de estudio:} & Periodo académico B26, con una duración de seis meses a partir de la aprobación de la propuesta (2026--2027) \\
    \hline
  \end{tabular}
  \raggedright
  \elaboradopor{Ramón Alejandro Manzaba Castro}
  \fuente{Elaboración propia}
\end{table}

La delimitación del área confirma que el trabajo pertenece a la gobernanza de la información y no solo a la seguridad técnica: su objeto es un proceso institucional con implicaciones jurídicas y no un sistema aislado.

\section{Metodología de Desarrollo o Implementación}

\subsection{Enfoque y tipo de investigación}

La investigación es de tipo aplicado, porque busca resolver un problema concreto de cumplimiento normativo en las IES, y de alcance descriptivo-propositivo, porque describe la situación de un proceso de tratamiento y propone un artefacto metodológico para gestionarlo. Se adopta un enfoque mixto: el diagnóstico y la validación combinan información cualitativa, obtenida mediante revisión documental, entrevistas y juicio de expertos, con información cuantitativa, como los niveles de madurez, las estimaciones de probabilidad y las métricas de calibración. El diseño es no experimental, dado que no se manipulan variables, sino que se observa y modela un proceso en su contexto \cite{hernandez2018metodologia}.

\subsection{Ciencia del diseño como marco metodológico}

Dado que el producto principal del trabajo es un artefacto ---una metodología acompañada de una matriz de riesgos reutilizable---, el desarrollo se rige por la investigación en ciencia del diseño (\textit{Design Science Research}, DSR). Hevner et al. \cite{hevner2004design} establecen que la ciencia del diseño en sistemas de información produce artefactos, como constructos, modelos, métodos e instanciaciones, que deben ser rigurosamente evaluados en función de su utilidad para resolver un problema relevante. Peffers et al. \cite{peffers2007dsrm} operacionalizaron este paradigma en un proceso de seis actividades: identificación del problema y motivación, definición de los objetivos de la solución, diseño y desarrollo, demostración, evaluación y comunicación. La Tabla~\ref{tab:dsr_fases} relaciona esas actividades con las fases de la investigación definidas en la propuesta, los objetivos específicos y los entregables.

\begin{table}[htbp]
  \centering
  \caption{Relación entre las actividades DSR, las fases de la investigación y los objetivos específicos.}
  \label{tab:dsr_fases}
  \footnotesize
  \renewcommand{\arraystretch}{1.2}
  \begin{tabularx}{\textwidth}{|p{3.1cm}|p{2.7cm}|c|L|}
    \hline
    \textbf{Actividad DSR} & \textbf{Fase de la investigación} & \textbf{OE} & \textbf{Entregable} \\
    \hline
    1. Identificación del problema y motivación & Diagnóstico situacional & OE1, OE2 & Cifras del sector (Sección~\ref{sec:cifras_sector}); requisitos del diagnóstico; caracterización del proceso, RAT preliminar e informe de madurez \\
    \hline
    2. Definición de los objetivos de la solución & Diagnóstico situacional & OE2 & Situación de las universidades, referentes comparados y requisitos de la metodología (Sección~\ref{sec:requerimientos}) \\
    \hline
    3. Diseño y desarrollo & Diseño metodológico & OE2 & Metodología por fases, criterios de estimación y plantilla de matriz (Sección~\ref{sec:diseno_solucion}) \\
    \hline
    4. Demostración & Aplicación en el caso de estudio & OE3 & Expediente anonimizado del caso, matriz de riesgos inherentes y residuales y plan de tratamiento \\
    \hline
    5. Evaluación & Validación metodológica & OE4 & Calibración, rondas Delphi, juicio de expertos, indicadores con metas y recomendaciones de mejora continua \\
    \hline
    6. Comunicación & Documentación final & OE4 & Trabajo de integración curricular y anexos \\
    \hline
  \end{tabularx}
  \raggedright
  \elaboradopor{Ramón Alejandro Manzaba Castro}
  \fuente{Elaboración propia con base en \cite{peffers2007dsrm}}
\end{table}

\subsection{Unidad de análisis, población y muestra}

La unidad de análisis es el proceso de admisión y matrícula de la UTI, sede Quito, desagregado en los seis subprocesos de la Tabla~\ref{tab:subprocesos}. El estudio no recolecta datos personales de estudiantes ni de aspirantes: la información se refiere a procesos, sistemas, controles y volúmenes agregados.

Se distinguen dos grupos de informantes. El primero está formado por los responsables de las unidades que intervienen en el proceso ---Admisiones, Secretaría Académica, Bienestar Estudiantil, Dirección Financiera y Dirección de Tecnologías de la Información, o sus equivalentes en la estructura de la UTI---, seleccionados de forma intencional por su conocimiento directo del subproceso a su cargo. El segundo es un panel de expertos, cuya población está constituida por profesionales con experiencia en protección de datos personales, seguridad de la información, gestión académica o derecho informático. La muestra es no probabilística e intencional \cite{hernandez2018metodologia}, compuesta por un panel de al menos cinco expertos, número mínimo que la guía de la SPDP establece para la estimación mediante el método Delphi \cite{spdp2026guiav2}. Los criterios de selección son: (a) título de tercer nivel en tecnologías de la información, sistemas, derecho o afines; (b) al menos tres años de experiencia en seguridad de la información, protección de datos o gestión de procesos académicos en IES, y (c) ausencia de conflicto de interés con el caso. Los expertos se identifican mediante códigos (E01, E02, \ldots) y participan con consentimiento informado.

En suma, el diseño metodológico responde a la naturaleza del problema: un artefacto que debe ser útil se construye y evalúa con ciencia del diseño, y un juicio sobre riesgos que carece de datos históricos se obtiene de expertos seleccionados con criterios explícitos.

\section{Técnicas e Instrumentos de Recolección de Información}

La Tabla~\ref{tab:tecnicas} resume las técnicas e instrumentos previstos, los informantes o fuentes y el objetivo al que tributa cada uno. Los instrumentos se presentan en los Anexos~\ref{anx:lista_verificacion} a~\ref{anx:consentimiento}, y la plantilla de la matriz de riesgos, en el Anexo~\ref{anx:plantilla_matriz}.

\begin{table}[htbp]
  \centering
  \caption{Técnicas e instrumentos de recolección de información.}
  \label{tab:tecnicas}
  \footnotesize
  \renewcommand{\arraystretch}{1.2}
  \begin{tabularx}{\textwidth}{|p{2.6cm}|L|p{3.0cm}|c|}
    \hline
    \textbf{Técnica} & \textbf{Instrumento} & \textbf{Fuente o informante} & \textbf{OE} \\
    \hline
    Revisión documental & Lista de verificación de requisitos de la LOPDP, su reglamento y la normativa de la SPDP aplicables al proceso & Normativa vigente; documentos institucionales autorizados (políticas, formularios, contratos) & OE1 \\
    \hline
    Entrevista semiestructurada & Guía de entrevista por subproceso (finalidades, datos, accesos, conservación, terceros, incidentes) & Responsables de las unidades del proceso & OE1 \\
    \hline
    Encuesta & Cuestionario de madurez en protección de datos, con escala de cinco niveles por dominio & Responsables de las unidades; validación de contenido por expertos & OE1 \\
    \hline
    Auditoría de procesos & Lista de verificación jurídica, organizacional y técnica orientada a la práctica efectiva y no solo a la política escrita & Evidencia del proceso recopilada con autorización & OE1, OE3 \\
    \hline
    Análisis documental comparado & Matriz de fuentes oficiales y estudios sobre universidades; tabla comparativa de regímenes & Registro Oficial, SPDP, CES, autoridades extranjeras, artículos indexados & OE2 \\
    \hline
    Autoevaluación de competencia & Cuestionario de coeficiente de competencia experta ($K$) & Candidatos al panel & OE4 \\
    \hline
    Método Delphi & Cuestionarios iterativos de estimación de probabilidad e impacto con justificación (\textit{rationale}) & Panel de al menos cinco expertos & OE3, OE4 \\
    \hline
    Juicio de expertos & Instrumento de validación de la metodología (pertinencia, claridad, suficiencia y coherencia), escala de Likert de cinco puntos & Panel de expertos & OE4 \\
    \hline
    Simulacro de vulneración & Protocolo de ejercicio de mesa con tiempos de notificación & Responsables del proceso y de TI & OE4 \\
    \hline
  \end{tabularx}
  \raggedright
  \elaboradopor{Ramón Alejandro Manzaba Castro}
  \fuente{Elaboración propia}
\end{table}

\subsection{Método Delphi y calibración de las estimaciones}

El método Delphi es una técnica estructurada de consulta a expertos que se desarrolla en rondas sucesivas, de forma anónima y con retroalimentación controlada, hasta alcanzar un nivel aceptable de consenso \cite{okoli2004delphi}. En este trabajo se emplea para estimar la probabilidad anual y el impacto de cada escenario de riesgo cuando no existen datos históricos suficientes, que es la situación habitual en las IES que inician su gestión de riesgos. Cada experto registra su estimación y su justificación; tras cada ronda se le devuelve la mediana y el rango intercuartílico del panel, junto con los argumentos anonimizados, y se le permite revisar su respuesta. Se prevén como máximo tres rondas.

\textbf{Selección por competencia.} Además de los criterios de perfil, cada candidato responde una autoevaluación (Anexo~\ref{anx:competencia}) con la que se calcula su coeficiente de competencia experta \cite{cabero2013competencia}, según la Ecuación~\ref{eq:competencia}:

\begin{equation}
\label{eq:competencia}
K = \tfrac{1}{2}\left(K_c + K_a\right)
\end{equation}
\addeqtolist{Coeficiente de competencia experta}

\noindent donde $K_c$ es el coeficiente de conocimiento, igual a la autovaloración del experto en una escala de 0 a 10 multiplicada por 0,1, y $K_a$ es el coeficiente de argumentación, que suma los pesos de las fuentes que sustentan su criterio (análisis teóricos, experiencia práctica, trabajos nacionales e internacionales, conocimiento del estado del tema e intuición). La competencia se considera alta si $K > 0{,}8$, media si $0{,}5 \le K \le 0{,}8$ y baja si $K < 0{,}5$ \cite{cabero2013competencia}. Se admiten en el panel los candidatos de competencia alta y, solo si no se alcanza el mínimo de cinco, los de competencia media, lo que se declara en los resultados.

\textbf{Criterio de consenso.} Para cada escenario y variable se considera alcanzado el consenso cuando se cumplen a la vez dos condiciones: el rango intercuartílico es menor o igual que 1 y al menos el 70\,\% de los expertos se sitúa a no más de un nivel de la mediana. Como medida complementaria del acuerdo del panel sobre el orden de prioridad de los escenarios, se calcula en cada ronda el coeficiente de concordancia $W$ de Kendall \cite{kendall1939rankings}, cuya significación se contrasta con el estadístico $\chi^2 = m(n-1)W$, con $n-1$ grados de libertad, siendo $m$ el número de expertos y $n$ el número de escenarios. El umbral del 70\,\% es una decisión de diseño de esta metodología, que el panel puede revisar en la primera ronda.

Para cumplir la exigencia de la guía de la SPDP de validar las estimaciones \cite{spdp2026guiav2}, los expertos responden además una batería de preguntas de calibración con respuesta conocida, expresando su grado de confianza como una probabilidad. La calidad de esas respuestas se mide con el puntaje de Brier \cite{brier1950verification}, definido en la Ecuación~\ref{eq:brier}:

\begin{equation}
\label{eq:brier}
BS = \frac{1}{N} \sum_{t=1}^{N} \left( f_t - o_t \right)^2
\end{equation}
\addeqtolist{Puntaje de Brier para la calibración de estimaciones}

\noindent donde $N$ es el número de preguntas, $f_t$ es la probabilidad asignada por el experto a la pregunta $t$ y $o_t$ toma el valor 1 si el evento ocurrió (o la afirmación es verdadera) y 0 en caso contrario. El puntaje varía entre 0, que corresponde a una calibración perfecta, y 1; un experto que respondiera siempre con una probabilidad de 0,5 obtendría 0,25, valor que se adopta como umbral de referencia. Las estimaciones de los expertos calibrados se emplean para ponderar las respuestas del panel conforme a la Ecuación~\ref{eq:ponderacion}, que es la que implementa la hoja \textit{Delphi} de la plantilla:

\begin{equation}
\label{eq:ponderacion}
\hat{x} = \frac{\sum_{e=1}^{m} w_e\, x_e}{\sum_{e=1}^{m} w_e}, \qquad
w_e = \begin{cases} 1 - BS_e & \text{si } BS_e \le 0{,}25 \\ 0 & \text{si } BS_e > 0{,}25 \end{cases}
\end{equation}
\addeqtolist{Estimación del panel ponderada por calibración}

\noindent donde $x_e$ es la estimación ordinal (1 a 5) del experto $e$ y $BS_e$ su puntaje de Brier. La mediana sigue siendo el valor de referencia del consenso; la media ponderada se informa como contraste, y si ambas difieren en más de un nivel, el escenario vuelve a ronda.

\subsection{Juicio de expertos sobre la metodología}

La validez de contenido de la metodología se valora con el instrumento del Anexo~\ref{anx:validacion}, que agrupa sus enunciados en cuatro criterios ---pertinencia, claridad, suficiencia y coherencia--- calificados en una escala de 1 a 5. Para cada enunciado se calcula la V de Aiken \cite{aiken1985three}, según la Ecuación~\ref{eq:aiken}:

\begin{equation}
\label{eq:aiken}
V = \frac{S}{n\,(c-1)}, \qquad S = \sum_{i=1}^{n} (x_i - l)
\end{equation}
\addeqtolist{Coeficiente V de Aiken}

\noindent donde $x_i$ es la calificación del juez $i$, $l = 1$ es el valor mínimo de la escala, $c = 5$ es el número de categorías y $n$ el número de jueces. Un enunciado se considera validado cuando $V \ge 0{,}70$ y la mediana del panel es igual o mayor que 4; un criterio se considera validado cuando lo están todos sus enunciados. El umbral de 0,70 es una decisión de diseño que se declara para que el resultado sea verificable; la significación estadística de cada $V$ se informa además con las tablas de probabilidad del propio Aiken \cite{aiken1985three}.

\subsection{Consideraciones éticas}

El trabajo se rige por las siguientes reglas: (a) la recolección de información en la UTI requiere la autorización formal de la universidad, a través de la instancia que esta determine, que se incorpora como Anexo~\ref{anx:autorizacion} y sin la cual no se aplica ningún instrumento; (b) no se incorporan al documento ni al repositorio del proyecto datos personales reales: la evidencia se registra en forma anonimizada, seudonimizada o agregada; (c) los informantes y los expertos se identifican con códigos (E01, E02, \ldots) y firman un consentimiento informado que describe la finalidad, el uso de la información y su derecho a retirarse; (d) los hallazgos se redactan en términos de procesos y controles, nunca de personas, y (e) el autor declara que desempeña funciones de administración de sistemas en la infraestructura de la UTI; si alguno de los sistemas que administra forma parte del proceso, se declarará expresamente y solo se utilizará la información que la institución autorice, sin aprovechar el acceso privilegiado del autor.

Estas reglas no son un trámite accesorio: una investigación sobre protección de datos que no protegiera los datos de sus propios informantes contradiría su objeto. Por ello, la ética del estudio se trata como parte del método y no como un apartado formal.

\section{Requerimientos del Proyecto}
\label{sec:requerimientos}

Dado que el artefacto es una metodología y no un sistema informático, los requerimientos se expresan como requisitos que la metodología debe satisfacer. Los requisitos funcionales describen lo que la metodología debe permitir hacer y se derivan de la LOPDP y de la guía de la SPDP; los no funcionales describen las cualidades que debe tener para ser adoptada por una IES.

\subsection{Requerimientos Funcionales}

La Tabla~\ref{tab:rf} enumera los requisitos funcionales y su fundamento normativo.

\begin{table}[htbp]
  \centering
  \caption{Requisitos funcionales de la metodología.}
  \label{tab:rf}
  \footnotesize
  \renewcommand{\arraystretch}{1.2}
  \begin{tabularx}{\textwidth}{|c|L|p{3.6cm}|}
    \hline
    \textbf{ID} & \textbf{Requisito} & \textbf{Fundamento} \\
    \hline
    RF01 & Diagnosticar el nivel de madurez de la institución en protección de datos antes de iniciar la gestión de riesgos. & LOPDP art.~37; fase 0 propuesta \\
    \hline
    RF02 & Inventariar las actividades de tratamiento (RAT) con finalidades, categorías de datos, titulares, destinatarios, encargados y plazos de conservación. & LOPDP art.~40; guía SPDP (contexto) \\
    \hline
    RF03 & Identificar vulnerabilidades de cuatro tipos: jurídicas, organizacionales, técnicas y relativas a grupos vulnerables de titulares. & Guía SPDP \\
    \hline
    RF04 & Construir escenarios de riesgo con comunidad de amenaza, vector, vulnerabilidades e impacto primario y secundario, desde la perspectiva del titular y del responsable. & Guía SPDP \\
    \hline
    RF05 & Estimar la probabilidad referida a un lapso anual y el impacto, con justificación documentada de cada valor. & Guía SPDP \\
    \hline
    RF06 & Modelar por separado la confidencialidad, la integridad y la disponibilidad, además de la licitud del tratamiento. & Guía SPDP; LOPDP art.~37 \\
    \hline
    RF07 & Ponderar el impacto con los criterios de categorías especiales, grupos vulnerables, número de titulares, naturaleza de la vulneración y volumen de datos. & LOPDP arts.~25 y 40; guía SPDP \\
    \hline
    RF08 & Determinar el nivel de riesgo inherente y residual con una regla explícita para los escenarios de baja probabilidad y muy alto impacto. & Guía SPDP \\
    \hline
    RF09 & Determinar la necesidad de una evaluación de impacto (EIPD) a partir de las etapas previas. & LOPDP art.~42; Res. SPDP-SPD-2026-0005-R \\
    \hline
    RF10 & Seleccionar opciones de tratamiento (mitigar, transferir, evitar, aceptar) y controles con responsable, plazo e indicador. & LOPDP art.~37; ISO/IEC 27005 \\
    \hline
    RF11 & Definir indicadores de monitoreo y un procedimiento de respuesta que permita notificar vulneraciones en el término legal. & LOPDP arts.~37 y 43 \\
    \hline
  \end{tabularx}
  \raggedright
  \elaboradopor{Ramón Alejandro Manzaba Castro}
  \fuente{Elaboración propia con base en \cite{lopdp2021,spdp2026guiav2,spdp2026granescala,iso27005_2022}}
\end{table}

\subsection{Requerimientos No Funcionales}

La Tabla~\ref{tab:rnf} presenta los requisitos no funcionales.

\begin{table}[htbp]
  \centering
  \caption{Requisitos no funcionales de la metodología.}
  \label{tab:rnf}
  \footnotesize
  \renewcommand{\arraystretch}{1.2}
  \begin{tabularx}{\textwidth}{|c|p{2.8cm}|L|}
    \hline
    \textbf{ID} & \textbf{Atributo} & \textbf{Descripción} \\
    \hline
    RNF01 & Trazabilidad & Cada actividad y cada campo de la matriz se vinculan con un objetivo específico, una etapa de la guía de la SPDP y un artículo de la LOPDP. \\
    \hline
    RNF02 & Reutilización & La matriz se entrega como plantilla de hoja de cálculo con listas desplegables, fórmulas y formato condicional, aplicable a otros procesos de la institución. \\
    \hline
    RNF03 & Usabilidad & La metodología puede ser aplicada por personal de TI, jurídico o de gestión sin formación especializada en análisis cuantitativo de riesgos. \\
    \hline
    RNF04 & Escalabilidad & Puede aplicarse a un proceso, a una unidad o a toda la institución sin cambiar su estructura. \\
    \hline
    RNF05 & Bajo costo & Se apoya en normas de acceso público y en herramientas ofimáticas de uso común. \\
    \hline
    RNF06 & Confidencialidad & Los expedientes de aplicación usan datos seudonimizados o agregados y registran el acceso a la evidencia. \\
    \hline
    RNF07 & Mantenibilidad & Las escalas, catálogos y criterios se mantienen en hojas separadas para actualizarlos ante cambios normativos. \\
    \hline
  \end{tabularx}
  \raggedright
  \elaboradopor{Ramón Alejandro Manzaba Castro}
  \fuente{Elaboración propia}
\end{table}

Los requisitos cumplen una función argumentativa precisa: fijan, antes del diseño, el criterio con el que se juzgará si la metodología es adecuada. Así, la evaluación del Capítulo~\ref{chap:implementacion} no depende de la opinión del autor, sino del cumplimiento verificable de cada requisito.

\section{Diseño de la Solución}
\label{sec:diseno_solucion}

\subsection{Insumos del análisis de universidades y referentes}
\label{sec:insumos_oe2}

El segundo objetivo específico no termina en un diagnóstico del sector, sino en decisiones de diseño. La Tabla~\ref{tab:insumos_oe2} muestra cómo cada hallazgo del análisis de las universidades y de los referentes nacionales e internacionales (Capítulo~\ref{chap:introduccion}) se traduce en un componente concreto de la metodología.

\begin{table}[htbp]
  \centering
  \caption{De los hallazgos sobre universidades y referentes a las decisiones de diseño (OE2).}
  \label{tab:insumos_oe2}
  \scriptsize
  \renewcommand{\arraystretch}{1.2}
  \begin{tabularx}{\textwidth}{|L|L|>{\raggedright\arraybackslash}p{2.6cm}|}
    \hline
    \textbf{Hallazgo} & \textbf{Decisión de diseño} & \textbf{Componente} \\
    \hline
    La adopción de políticas de protección de datos en las universidades es desigual y no existe una medición oficial de su implementación \cite{piray2026garantias} & Diagnosticar la madurez antes de gestionar el riesgo y verificar la práctica efectiva, no solo la política escrita & Fase 0; Anexos~\ref{anx:lista_verificacion} y~\ref{anx:madurez} \\
    \hline
    Los delegados registrados equivalían como máximo al 7,0\,\% de las instituciones educativas (Sección~\ref{sec:cifras_sector}) & La metodología debe poder aplicarse sin un especialista dedicado y asignar las observaciones al área jurídica cuando no exista delegado & RNF03; regla RA-5 \\
    \hline
    El tratamiento de datos de menores en entornos educativos es gran escala por disposición directa \cite{spdp2026granescala} & Incorporar el Modelo Técnico de Gran Escala como agravante y como disparador de la EIPD & Agravante (b); RF09; regla RA-4 \\
    \hline
    Los plazos de notificación son de 5 días a la autoridad y 3 días a los titulares, y 72 horas para incidentes de ciberseguridad en entidades públicas (Sección~\ref{sec:regimen_vulneraciones}) & Preparar el procedimiento de respuesta antes de evaluar y comprobarlo con un simulacro & Fase P; fase 4; escenario de validación V2 \\
    \hline
    El compromiso de la alta dirección es un factor crítico de éxito en universidades \cite{fernandes2022csf} & Exigir autorización y patrocinio de la máxima autoridad como primera actividad & Fase P \\
    \hline
    Ningún marco revisado combina enfoque en el titular, ciclo completo, EIPD, estimación justificada, sector IES y normativa ecuatoriana (Tabla~\ref{tab:comparativa_marcos}) & Integrar la guía de la SPDP con ISO/IEC 27005, ISO/IEC 27035 y NIST SP 800-61r3, con escalas de impacto centradas en el titular & Estructura de fases; Tabla~\ref{tab:equivalencia_marcos} \\
    \hline
    España exige delegado a todas las universidades; Chile lo mantiene voluntario; Ecuador lo exige en el sector público y en gran escala (Tabla~\ref{tab:comparativa_paises}) & Tratar la existencia del delegado como dato del diagnóstico y no como supuesto & Requisito D09 \\
    \hline
  \end{tabularx}
  \raggedright
  \elaboradopor{Ramón Alejandro Manzaba Castro}
  \fuente{Elaboración propia}
\end{table}

\subsection{Estructura general de la metodología}

La metodología propuesta se compone de una fase de preparación (fase P), una fase de diagnóstico de madurez (fase 0) y cuatro fases de gestión: identificación, evaluación, mitigación y monitoreo. La fase de preparación se incorporó porque tanto la gestión de incidentes de ISO/IEC 27035-1:2023 como la guía de respuesta de NIST SP 800-61r3 sitúan antes de cualquier evaluación las condiciones que la hacen posible: compromiso de la dirección, roles, criterios, procedimientos y capacidades \cite{iso27035_1_2023,iso27035_2_2023,nist2025sp80061r3}. Sin ellas, la institución no podría, por ejemplo, notificar una vulneración dentro del término legal. Las fases se articulan con las etapas obligatorias de la guía de la SPDP, que a su vez sigue el proceso de ISO/IEC 27005:2022 \cite{spdp2026guiav2,iso27005_2022}, y se apoyan en dos actividades transversales: la comunicación y consulta con las partes interesadas, y el registro documental de cada decisión y de su justificación. La Figura~\ref{fig:fases_metodologia} representa esta estructura.

\begin{figure}[htbp]
  \centering
  \begin{tikzpicture}[
      font=\scriptsize,
      fase/.style={draw, rounded corners=3pt, fill=gray!12, align=center, text width=1.85cm, minimum height=1.55cm, inner sep=3pt},
      trans/.style={draw, align=center, text width=13.9cm, minimum height=0.6cm, inner sep=2pt},
      flecha/.style={-{Stealth[length=2mm]}, thick}
    ]
    \node[fase] (FP) at (0,0)      {\textbf{Fase P}\\Preparación\\(autorización, roles, criterios)};
    \node[fase] (F0) at (2.4,0)    {\textbf{Fase 0}\\Diagnóstico de madurez\\(OE1)};
    \node[fase] (F1) at (4.8,0)    {\textbf{Fase 1}\\Identificación\\(contexto e identificación)};
    \node[fase] (F2) at (7.2,0)    {\textbf{Fase 2}\\Evaluación\\(análisis, evaluación y EIPD)};
    \node[fase] (F3) at (9.6,0)    {\textbf{Fase 3}\\Mitigación\\(tratamiento)};
    \node[fase] (F4) at (12.0,0)   {\textbf{Fase 4}\\Monitoreo\\(seguimiento y respuesta)};
    \foreach \a/\b in {FP/F0,F0/F1,F1/F2,F2/F3,F3/F4} { \draw[flecha] (\a.east) -- (\b.west); }
    \draw[flecha] (F4.north) -- ++(0,0.55) -| node[pos=0.25, above] {Revisión periódica y ante cambios o incidentes} (F1.north);
    \node[trans] at (6.0,-1.45) {Comunicación y consulta con partes interesadas (DPD, TI, jurídico, unidades académicas)};
    \node[trans] at (6.0,-2.2)  {Registro documental y justificación (\textit{rationale}) de cada estimación y decisión};
  \end{tikzpicture}
  \caption{Estructura de la metodología propuesta.}
  \label{fig:fases_metodologia}
  \raggedright
  \elaboradopor{Ramón Alejandro Manzaba Castro}
  \fuente{Elaboración propia con base en \cite{spdp2026guiav2,iso27005_2022}}
\end{figure}

La Tabla~\ref{tab:fases_detalle} detalla, para cada fase, las actividades principales, los productos, la etapa correspondiente de la guía de la SPDP y el fundamento en la LOPDP. Este detalle se aplica al proceso de admisión y matrícula de la UTI en el Capítulo~\ref{chap:implementacion}.

\begin{table}[htbp]
  \centering
  \caption{Fases, actividades y productos de la metodología.}
  \label{tab:fases_detalle}
  \footnotesize
  \renewcommand{\arraystretch}{1.2}
  \begin{tabularx}{\textwidth}{|>{\raggedright\arraybackslash}p{1.9cm}|L|>{\raggedright\arraybackslash}p{3.0cm}|>{\raggedright\arraybackslash}p{2.3cm}|}
    \hline
    \textbf{Fase} & \textbf{Actividades principales} & \textbf{Productos} & \textbf{Etapa SPDP / LOPDP} \\
    \hline
    P. Preparación & Obtener la autorización y el patrocinio de la máxima autoridad; definir alcance, roles (responsable de proceso, DPD, TI, jurídico) y criterios de aceptación; conformar y calibrar el panel; establecer el procedimiento de respuesta a vulneraciones y los canales de notificación (SISPDP, organismo de control, CSIRT). & Autorización; plan de trabajo con roles; panel calibrado; procedimiento de respuesta aprobado & Contexto (criterios) / arts.~37, 43, 46 y 47 \\
    \hline
    0. Diagnóstico de madurez & Definir alcance; aplicar el cuestionario de madurez; revisar documentos y práctica efectiva; entrevistar a responsables. & Informe de madurez; RAT preliminar; brechas prioritarias & Previa al contexto / art.~37 \\
    \hline
    1. Identificación & Establecer el contexto interno, externo y normativo; completar el RAT; identificar activos, amenazas y vulnerabilidades de los cuatro tipos; construir escenarios de riesgo. & RAT; catálogo de amenazas y vulnerabilidades; escenarios de riesgo & Contexto e identificación / arts.~25 y 40 \\
    \hline
    2. Evaluación & Estimar probabilidad anual e impacto con justificación (Delphi); determinar el nivel inherente; contrastar con los criterios de aceptación; decidir si se requiere EIPD. & Matriz de riesgo inherente; mapa de calor; dictamen de EIPD & Análisis y evaluación / arts.~40 y 42 \\
    \hline
    3. Mitigación & Elegir la opción de tratamiento; seleccionar controles jurídicos, organizacionales y técnicos; estimar el riesgo residual; aprobar el riesgo residual. & Plan de tratamiento; matriz de riesgo residual & Tratamiento / art.~37 \\
    \hline
    4. Monitoreo & Definir indicadores; revisar periódicamente los escenarios; gestionar y notificar vulneraciones; actualizar la matriz y retroalimentar la fase 1. & Tablero de indicadores; procedimiento de respuesta; registro de revisiones & Monitoreo y comunicación / arts.~37 y 43 \\
    \hline
  \end{tabularx}
  \raggedright
  \elaboradopor{Ramón Alejandro Manzaba Castro}
  \fuente{Elaboración propia con base en \cite{lopdp2021,spdp2026guiav2,iso27005_2022}}
\end{table}

\subsection{Equivalencia con marcos de gestión de riesgos e incidentes}
\label{sec:equivalencia_marcos}

La Tabla~\ref{tab:equivalencia_marcos} muestra la correspondencia de cada fase con las etapas del proceso de gestión del riesgo de ISO 31000 e ISO/IEC 27005, que adopta la guía de la SPDP; con las cinco fases de gestión de incidentes de ISO/IEC 27035-1:2023 ---planificar y preparar, detectar y reportar, evaluar y decidir, responder y aprender lecciones---, y con las funciones del CSF 2.0 en que NIST SP 800-61r3 organiza la respuesta a incidentes. Esta última revisión, publicada en abril de 2025, sustituyó el ciclo de cuatro fases de la revisión 2 por las funciones Gobernar, Identificar, Proteger, Detectar, Responder y Recuperar, y agrupa las tres primeras como preparación y las tres últimas como respuesta \cite{nist2025sp80061r3}.

\begin{table}[htbp]
  \centering
  \caption{Equivalencia de las fases de la metodología con ISO 31000/ISO/IEC 27005, ISO/IEC 27035-1 y NIST SP 800-61r3.}
  \label{tab:equivalencia_marcos}
  \scriptsize
  \renewcommand{\arraystretch}{1.2}
  \begin{tabularx}{\textwidth}{|>{\raggedright\arraybackslash}p{2.3cm}|L|L|L|}
    \hline
    \textbf{Fase de la metodología} & \textbf{Guía SPDP / ISO 31000 / ISO/IEC 27005:2022} & \textbf{ISO/IEC 27035-1:2023} & \textbf{NIST SP 800-61r3 (funciones CSF 2.0)} \\
    \hline
    P. Preparación & Alcance, contexto y criterios; comunicación y consulta & Planificar y preparar & Gobernar (GV); Proteger (PR) en lo relativo a capacidades de respuesta \\
    \hline
    0. Diagnóstico de madurez & Establecimiento del contexto (interno) & Planificar y preparar (evaluación de capacidades) & Identificar (ID), incluida la mejora (ID.IM) \\
    \hline
    1. Identificación & Contexto e identificación del riesgo & --- (insumo para detectar y reportar) & Identificar (ID.AM, ID.RA) \\
    \hline
    2. Evaluación & Análisis y evaluación del riesgo; EIPD & Evaluar y decidir (criterios de severidad) & Identificar (ID.RA) \\
    \hline
    3. Mitigación & Tratamiento del riesgo & --- (controles preventivos) & Proteger (PR) \\
    \hline
    4. Monitoreo y respuesta & Seguimiento y revisión; registro e informe & Detectar y reportar; evaluar y decidir; responder; aprender lecciones & Detectar (DE); Responder (RS); Recuperar (RC); Identificar, mejora (ID.IM) \\
    \hline
  \end{tabularx}
  \raggedright
  \elaboradopor{Ramón Alejandro Manzaba Castro}
  \fuente{Elaboración propia con base en \cite{iso31000_2018,iso27005_2022,spdp2026guiav2,iso27035_1_2023,iso27035_2_2023,nist2025sp80061r3}}
\end{table}

La equivalencia muestra que la metodología no reemplaza a estos marcos, sino que los integra: de ISO/IEC 27005 toma el proceso de gestión del riesgo, de ISO/IEC 27035 y NIST SP 800-61r3 toma la preparación y la respuesta, y añade lo que ninguno de ellos aporta por sí solo, que es la valoración del daño a los titulares y los plazos de la normativa ecuatoriana.

\subsection{Modelo de estimación del riesgo}

El modelo de estimación respeta las reglas de la guía de la SPDP: cada escenario se analiza por separado para la confidencialidad, la integridad, la disponibilidad y la licitud; la probabilidad se refiere a un lapso anual; cada valor se acompaña de su justificación, y las escalas se tratan como etiquetas ordinales, no como magnitudes que puedan multiplicarse \cite{spdp2026guiav2}. Las escalas que se proponen a continuación constituyen la versión inicial que el panel Delphi revisará y calibrará.

La Tabla~\ref{tab:escala_probabilidad} define la escala de probabilidad. Cada nivel se asocia con un intervalo de frecuencia anual esperada, lo que obliga a que la estimación tenga un referente verificable y permite contrastarla en el futuro con el registro de incidentes.

\begin{table}[htbp]
  \centering
  \caption{Escala de probabilidad (lapso anual).}
  \label{tab:escala_probabilidad}
  \small
  \renewcommand{\arraystretch}{1.2}
  \begin{tabularx}{\textwidth}{|p{2.2cm}|p{3.6cm}|L|}
    \hline
    \textbf{Nivel} & \textbf{Frecuencia anual esperada} & \textbf{Descripción orientativa} \\
    \hline
    Muy baja & Menos de una vez cada 10 años & No hay antecedentes en la institución ni en el sector; la amenaza requiere capacidades poco comunes. \\
    \hline
    Baja & Una vez cada 2 a 10 años & Hay antecedentes en el sector, pero los controles existentes dificultan su materialización. \\
    \hline
    Media & Aproximadamente una vez cada 1 a 2 años & Hay antecedentes en IES similares y controles parciales. \\
    \hline
    Alta & Entre 1 y 5 veces por año & El evento ha ocurrido en la institución o la vulnerabilidad es conocida y explotable. \\
    \hline
    Muy alta & Más de 5 veces por año & El evento es recurrente o la vulnerabilidad se explota con facilidad. \\
    \hline
  \end{tabularx}
  \raggedright
  \elaboradopor{Ramón Alejandro Manzaba Castro}
  \fuente{Elaboración propia con base en \cite{spdp2026guiav2}}
\end{table}

El impacto se valora en primer lugar desde la perspectiva del titular, conforme al enfoque de derechos de la LOPDP, y en segundo lugar desde la del responsable (sanciones, costos y reputación). La Tabla~\ref{tab:escala_impacto} define la escala de impacto sobre el titular, adaptada de los niveles de severidad propuestos por ENISA \cite{enisa2013severity}.

\begin{table}[htbp]
  \centering
  \caption{Escala de impacto sobre el titular.}
  \label{tab:escala_impacto}
  \small
  \renewcommand{\arraystretch}{1.2}
  \begin{tabularx}{\textwidth}{|p{2.2cm}|L|}
    \hline
    \textbf{Nivel} & \textbf{Consecuencias para el titular} \\
    \hline
    Muy bajo & Molestias menores que el titular supera sin dificultad (por ejemplo, reingresar información). \\
    \hline
    Bajo & Inconvenientes que el titular supera con algún esfuerzo (costos menores, tiempo, molestias). \\
    \hline
    Medio & Consecuencias significativas que el titular puede superar con dificultad (pérdida de oportunidades académicas, cargos indebidos reversibles, temor o estrés). \\
    \hline
    Alto & Consecuencias graves, difíciles de superar (fraude o suplantación de identidad, discriminación, exclusión de beneficios, daño reputacional). \\
    \hline
    Muy alto & Consecuencias irreversibles o que afectan derechos fundamentales, la salud o la integridad del titular (exposición de datos de salud o discapacidad, afectación a menores de edad). \\
    \hline
  \end{tabularx}
  \raggedright
  \elaboradopor{Ramón Alejandro Manzaba Castro}
  \fuente{Elaboración propia adaptada de \cite{enisa2013severity}}
\end{table}

El nivel de impacto inicial se ajusta con los criterios agravantes que exige la guía de la SPDP \cite{spdp2026guiav2}: (a) si el escenario involucra categorías especiales de datos o datos de grupos vulnerables, en particular de adolescentes \cite[art.~25]{lopdp2021}, el impacto sube un nivel; (b) si el tratamiento afectado alcanza el umbral de gran escala del Modelo Técnico de Gran Escala de la SPDP o encaja en uno de sus supuestos directos \cite{spdp2026granescala}, o supera el umbral de titulares que la institución haya definido como masivo, el impacto sube un nivel; (c) si la vulneración es intencional o los datos pueden combinarse para identificar a los titulares con facilidad, el impacto sube un nivel. El ajuste acumulado no puede superar el nivel muy alto y cada ajuste debe justificarse.

\textbf{Fuentes de información de cada estimación.} La guía de la SPDP exige que todo valor tenga una justificación \cite{spdp2026guiav2}; para que esa justificación no dependa del criterio del analista, la Tabla~\ref{tab:fuentes_estimacion} fija de antemano qué fuente respalda cada variable y en qué orden de preferencia se usan.

\begin{table}[htbp]
  \centering
  \caption{Fuentes de información y reglas para cada variable de la estimación.}
  \label{tab:fuentes_estimacion}
  \scriptsize
  \renewcommand{\arraystretch}{1.2}
  \begin{tabularx}{\textwidth}{|>{\raggedright\arraybackslash}p{2.4cm}|L|L|}
    \hline
    \textbf{Variable} & \textbf{Fuentes, en orden de preferencia} & \textbf{Regla de asignación} \\
    \hline
    Probabilidad anual ($P$) & (1) Registro de incidentes de la institución (D08); (2) estado de los controles verificado en el diagnóstico (D06); (3) antecedentes sectoriales documentados \cite{sophos2024education,ulven2021systematic,li2023breach}; (4) estimación Delphi calibrada & Se asigna el nivel cuya frecuencia anual (Tabla~\ref{tab:escala_probabilidad}) contiene la frecuencia estimada; si las fuentes discrepan, prevalece la de mayor preferencia y se registra la discrepancia \\
    \hline
    Impacto sobre el titular ($I_0$) & (1) Categorías de datos y titulares del RAT (D03); (2) criterios de severidad de ENISA \cite{enisa2013severity}; (3) estimación Delphi calibrada & Se asigna el nivel cuyo descriptor (Tabla~\ref{tab:escala_impacto}) corresponde a la consecuencia más grave razonablemente previsible para el titular \\
    \hline
    Agravantes ($a_1, a_2, a_3$) & RAT (D03); volumen y proporción de menores (D04); Modelo Técnico de Gran Escala \cite{spdp2026granescala} & Cada agravante vale 0 o 1 según se cumpla o no el criterio (a), (b) o (c) \\
    \hline
    Impacto sobre el responsable & Régimen sancionatorio \cite[arts.~71 y 72]{lopdp2021}; volumen de negocio; costos de respuesta & Se registra como información para la decisión; no modifica el nivel, que se determina por el impacto sobre el titular \\
    \hline
  \end{tabularx}
  \raggedright
  \elaboradopor{Ramón Alejandro Manzaba Castro}
  \fuente{Elaboración propia con base en \cite{spdp2026guiav2,enisa2013severity,spdp2026granescala,lopdp2021}}
\end{table}

Formalmente, el impacto ajustado se obtiene con la Ecuación~\ref{eq:impacto_ajustado}, en la que los niveles se codifican de 1 (muy bajo) a 5 (muy alto) solo para aplicar el tope, no para operar con ellos como magnitudes:

\begin{equation}
\label{eq:impacto_ajustado}
I_{aj} = \min\left(5,\; I_0 + a_1 + a_2 + a_3\right), \qquad a_k \in \{0, 1\}
\end{equation}
\addeqtolist{Impacto ajustado por agravantes}

\subsection{Determinación del nivel de riesgo y criterios de aceptación}

Para evitar que el nivel de riesgo se obtenga multiplicando etiquetas, práctica que la guía de la SPDP desaconseja \cite{spdp2026guiav2}, el nivel no se calcula con una fórmula aritmética, sino con una función de decisión $N = T(P, I_{aj})$, donde $T$ es la Tabla~\ref{tab:matriz_nivel}. La tabla otorga mayor peso al impacto que a la probabilidad y contiene una regla explícita para los escenarios de baja probabilidad y alto impacto: todo escenario con impacto muy alto sobre el titular se clasifica, como mínimo, en nivel alto, con independencia de su probabilidad. El mapa de calor resultante se utiliza solo para representar los resultados, no para calibrarlos.

\begin{table}[htbp]
  \centering
  \caption{Tabla de decisión del nivel de riesgo.}
  \label{tab:matriz_nivel}
  \small
  \renewcommand{\arraystretch}{1.35}
  \begin{tabular}{|l|c|c|c|c|c|}
    \hline
    \textbf{Probabilidad / Impacto} & \textbf{Muy bajo} & \textbf{Bajo} & \textbf{Medio} & \textbf{Alto} & \textbf{Muy alto} \\
    \hline
    \textbf{Muy alta} & Medio & Alto & Alto & Muy alto & Muy alto \\
    \hline
    \textbf{Alta}     & Bajo & Medio & Alto & Muy alto & Muy alto \\
    \hline
    \textbf{Media}    & Bajo & Medio & Medio & Alto & Muy alto \\
    \hline
    \textbf{Baja}     & Muy bajo & Bajo & Medio & Alto & Alto \\
    \hline
    \textbf{Muy baja} & Muy bajo & Muy bajo & Bajo & Medio & Alto \\
    \hline
  \end{tabular}
  \par\vspace{2pt}
  \raggedright
  \elaboradopor{Ramón Alejandro Manzaba Castro}
  \fuente{Elaboración propia con base en \cite{spdp2026guiav2}}
\end{table}

Los criterios de aceptación del riesgo se expresan como reglas numeradas, de modo que cualquier revisor pueda verificar su aplicación:

\begin{itemize}
  \item \textbf{RA-1.} Nivel muy bajo o bajo: aceptable, con registro de la decisión y revisión al cierre de cada periodo de admisión.
  \item \textbf{RA-2.} Nivel medio: se trata dentro del plan anual, con un plazo máximo de doce meses, salvo que la dirección acepte formalmente el riesgo con su justificación.
  \item \textbf{RA-3.} Nivel alto o muy alto: no aceptable; requiere tratamiento prioritario, con plazo máximo de noventa días para iniciar los controles.
  \item \textbf{RA-4.} Nivel alto o muy alto en escenarios con categorías especiales o en tratamientos a gran escala: activa la evaluación de impacto \cite[art.~42]{lopdp2021}.
  \item \textbf{RA-5.} Impacto ajustado muy alto ($I_{aj} = 5$): ningún riesgo residual puede aceptarse sin las observaciones previas del delegado de protección de datos y del área jurídica y sin la aprobación de la máxima autoridad.
\end{itemize}

Los plazos de las reglas RA-2 y RA-3 son parámetros de diseño propuestos por esta metodología; la institución puede ajustarlos en la fase de preparación, siempre que lo documente.

\subsection{Estructura de la matriz de riesgos}

La matriz de riesgos se implementa como una hoja de cálculo con una fila por escenario y por dimensión (confidencialidad, integridad, disponibilidad o licitud). Sus campos se agrupan en seis bloques, que se resumen en la Tabla~\ref{tab:estructura_matriz}. Hojas adicionales contienen las escalas de probabilidad e impacto, la tabla de decisión, los criterios de aceptación, el catálogo de amenazas, el catálogo de controles referidos a ISO/IEC 27002:2022 \cite{iso27002_2022}, los resultados del panel Delphi y el cálculo del puntaje de Brier.

\begin{table}[htbp]
  \centering
  \caption{Bloques de campos de la matriz de riesgos.}
  \label{tab:estructura_matriz}
  \footnotesize
  \renewcommand{\arraystretch}{1.2}
  \begin{tabularx}{\textwidth}{|p{2.8cm}|L|}
    \hline
    \textbf{Bloque} & \textbf{Campos} \\
    \hline
    Identificación del tratamiento & ID, subproceso, actividad de tratamiento, categorías de datos, categoría especial (sí/no), tipo y número estimado de titulares, grupo vulnerable, activo o sistema \\
    \hline
    Escenario de riesgo & Dimensión (C/I/D/licitud), comunidad de amenaza, vector, vulnerabilidad y su tipo (jurídica, organizacional, técnica, titulares), artículo de la LOPDP o norma de la SPDP, control existente \\
    \hline
    Riesgo inherente & Probabilidad anual y su justificación; impacto sobre el titular, impacto sobre el responsable y su justificación; agravantes aplicados; nivel inherente \\
    \hline
    Tratamiento & Opción (mitigar, transferir, evitar, aceptar), control propuesto con referencia a ISO/IEC 27002, responsable, plazo \\
    \hline
    Riesgo residual & Probabilidad residual, impacto residual, nivel residual, aceptación \\
    \hline
    Monitoreo & Indicador, frecuencia de revisión, estado, fecha de la última revisión \\
    \hline
  \end{tabularx}
  \raggedright
  \elaboradopor{Ramón Alejandro Manzaba Castro}
  \fuente{Elaboración propia}
\end{table}

El diseño, en conjunto, responde a cada requisito funcional con una actividad, un criterio o un campo de la matriz. Su rasgo distintivo frente a los modelos revisados en el Capítulo~\ref{chap:introduccion} es que el titular ocupa el centro de la valoración y que ninguna estimación queda sin justificación documentada.

\section{Justificación Técnica}

La selección de los componentes de la metodología responde a cinco decisiones.

\textbf{Marco de proceso.} Se adoptan las etapas de la guía de la SPDP, alineadas con ISO/IEC 27005:2022, porque la guía es la referencia normativa vigente en Ecuador y su capítulo~I es de cumplimiento obligatorio \cite{spdp2025guia,spdp2026guiav2}; ISO/IEC 27005 aporta el detalle operativo del proceso y facilita la integración con un sistema de gestión de seguridad de la información \cite{iso27005_2022}. Se descartó adoptar un marco extranjero completo, como el de la AEPD o el NIST Privacy Framework, porque ninguno responde a la LOPDP y su adaptación exigiría el mismo esfuerzo que construir la metodología sobre la guía nacional \cite{aepd2021riesgo,nist2020privacy}.

\textbf{Enfoque centrado en el titular.} A diferencia de los modelos aplicados en el sector público y educativo ecuatoriano, centrados en activos \cite{patino2018ict,lopezvasco2025iso}, la escala de impacto principal se refiere al daño al titular. La literatura reciente respalda incorporar categorías de daño a la privacidad en las evaluaciones de riesgo \cite{wairimu2025harm}.

\textbf{Estimación cualitativa calibrada.} Los modelos cuantitativos como FAIR requieren datos de frecuencia y pérdidas que una IES rara vez posee \cite{wang2020fair}. Por ello se opta por escalas ordinales ancladas en frecuencias anuales, estimadas mediante un panel Delphi y validadas con el puntaje de Brier, opción que la propia guía de la SPDP admite cuando se documenta la justificación de cada valor \cite{spdp2026guiav2}. La estructura de la matriz permite migrar a estimaciones cuantitativas cuando la institución acumule datos históricos.

\textbf{Herramienta de soporte.} La matriz se implementa en una hoja de cálculo, compatible con Microsoft Excel y LibreOffice Calc, en lugar de una plataforma comercial de gobierno, riesgo y cumplimiento. Esta decisión reduce el costo de adopción, facilita la revisión por parte de unidades no técnicas y cumple los requisitos de reutilización y mantenibilidad (RNF02, RNF05 y RNF07).

\textbf{Marco de investigación.} La ciencia del diseño es el paradigma adecuado porque el resultado es un artefacto cuya utilidad debe demostrarse y evaluarse \cite{hevner2004design}; el proceso de Peffers et al. \cite{peffers2007dsrm} ofrece una secuencia verificable que coincide con las fases de la investigación declaradas en la propuesta.

Podría sostenerse que una metodología cualitativa es menos rigurosa que un modelo cuantitativo. Sin embargo, el rigor no depende del tipo de escala, sino de que cada valor esté justificado, referido a un lapso y contrastado con evidencia, tres condiciones que el diseño propuesto exige. En conclusión, las decisiones adoptadas producen una metodología que cumple la normativa ecuatoriana, es aplicable con los recursos de una IES y deja abierta la evolución hacia estimaciones cuantitativas a medida que la institución acumule datos. El capítulo siguiente pone a prueba esa afirmación en el proceso de admisión y matrícula de la UTI.
